\section{Post-training III: RLVR y recompensas verificables}

\subsection{Por qué RLVR se convirtió en el punto de inflexión clave para la mejora del razonamiento}
El expositor define RLVR (Reinforcement Learning with Verifiable Rewards) como el punto de inflexión clave dentro del post-training: cuando una tarea puede verificarse de manera automática como correcta o incorrecta, el modelo puede mejorar de forma continua a partir de una gran cantidad de muestras autogeneradas, sin depender ya por completo de anotaciones costosas de preferencia humana.

\begin{importantbox}{El ciclo mínimo de RLVR}
\begin{enumerate}
    \item Generar respuestas candidatas;
    \item Dar una recompensa verificable mediante un verificador basado en reglas o en un programa;
    \item Actualizar la política con RL para aumentar la probabilidad de las muestras con alta recompensa.
\end{enumerate}
\end{importantbox}

\begin{figure}[H]
\centering
\includegraphics[width=0.9\textwidth]{frame-06-rlvr.jpg}
\caption{Video aproximadamente en 01:03:20: transición del reward model neuronal hacia recompensas verificables, lo que reduce la dependencia de anotaciones}
\end{figure}

\subsection{Comparación entre modelos de recompensa neuronales y recompensas basadas en reglas}
\begin{knowledgebox}{Por qué el expositor enfatiza «verificable»}
Los modelos de recompensa neuronales son flexibles, pero presentan sesgos y una interpretabilidad limitada; las recompensas basadas en reglas son más limitadas, pero más confiables. Para tareas como matemáticas, código y razonamiento simbólico, las recompensas basadas en reglas pueden reducir de manera significativa el riesgo de reward hacking.
\end{knowledgebox}

\begin{table}[H]
\centering
\renewcommand{\arraystretch}{1.2}
\begin{tabular}{p{2.8cm}p{3.6cm}p{3.8cm}p{2.8cm}}
\toprule
\textbf{Tipo de recompensa} & \textbf{Ventajas} & \textbf{Desventajas} & \textbf{Tarea típica} \\
\midrule
Modelo de recompensa neuronal & Alta capacidad expresiva, puede cubrir objetivos subjetivos complejos & Propenso a sesgos, interpretabilidad débil, costo alto & Preferencias de diálogo, consistencia de estilo \\
Recompensa basada en reglas/verificable & Estable, auditable, permite entrenamiento a escala automatizado & Cobertura limitada, costo de diseño alto & Matemáticas, código, restricciones formalizadas \\
\bottomrule
\end{tabular}
\caption{Comparación entre los dos tipos de mecanismos de recompensa}
\end{table}

\begin{warningbox}{Condiciones límite de RLVR}
RLVR no mejora de manera automática todas las tareas. Si una tarea carece de un verificador confiable, o si el verificador está desalineado con el objetivo real, el entrenamiento empujará al modelo hacia la dirección errónea de «alta puntuación, baja calidad».
\end{warningbox}

\subsection{De la experiencia de DeepSeek-R1 a prácticas reproducibles en la comunidad abierta}
El expositor cita la inspiración de la ruta al estilo R1: cuando el verificador es económico y las muestras pueden generarse a escala, el RL puede amplificar rápidamente la capacidad de razonamiento. Para la comunidad abierta, lo importante es liberar en conjunto el verificador, los scripts de entrenamiento y el proceso de evaluación como código abierto; de lo contrario, es difícil verificar si las conclusiones son transferibles.

\subsection{Resumen de la sección}
El verdadero valor de RLVR consiste en transformar la «mejora del razonamiento», de una alta dependencia manual, en una «iteración programable». La condición previa es que la calidad del verificador sea confiable y esté alineada con el objetivo real.

